\chapter{Teoretická část}

\section{Impedance}

Impedance $Z$ je fyzikální veličina, která popisuje celkový odpor elektrického obvodu vůči střídavému proudu o dané frekvenci.
Vyjadřuje se komplexním číslem, kdy reálná část odpovídá odporu $R$ a imaginární část reprezentuje reaktanci $X$.
\begin{equation}
Z = R + jX
\end{equation}

Impedance je často také vyjádřena v polárním tvaru:
\begin{equation}
Z = |Z| \cdot e^{j\varphi} = |Z| \angle \varphi
\end{equation}
kde $|Z| = \sqrt{R^2 + X^2}$ je modul impedance a $\varphi = \arctan(\frac{X}{R})$ je fázový úhel (argument).

V některých případech je výhodné použít obrácenou hodnotu impedance $Y = \frac{1}{Z}$, zvanou admitance. \cite{KeysightImpedanceHandbook}

\section{Veličiny odvozené z impedance}
Při měření impedance součástek a obvodů je často cílem určit jejich fyzikální parametry nebo charakteristiky, které jsou z impedance odvozeny.

\subsection{Primární parametry součástek}

Pro základní pasivní součástky lze z naměřené impedance určit jejich primární parametry:

\begin{itemize}
    \item \textbf{Odpor} $R$ -- reálná část impedance rezistoru
    \item \textbf{Indukčnost} $L$ -- vypočítá se z imaginární části impedance cívky podle vztahu:
    \begin{equation}
    L = \frac{X_L}{2\pi f} = \frac{\text{Im}(Z)}{2\pi f}
    \end{equation}
    \item \textbf{Kapacita} $C$ -- vypočítá se z imaginární části impedance kondenzátoru podle vztahu:
    \begin{equation}
    C = \frac{1}{2\pi f \cdot |X_C|} = \frac{-1}{2\pi f \cdot \text{Im}(Z)}
    \end{equation}
\end{itemize}

kde $f$ je frekvence měřicího signálu.

\subsection{Činitel kvality}

Činitel kvality $Q$ (quality factor) vyjadřuje poměr mezi reaktivní a rezistivní složkou impedance. \cite{KeysightImpedanceHandbook} Pro cívky a kondenzátory se definuje jako:
\begin{equation}
Q = \frac{|X|}{R} = \frac{|\text{Im}(Z)|}{\text{Re}(Z)}
\end{equation}

Vysoká hodnota $Q$ znamená nízké ztráty v součástce. Typicky se $Q$ používá pro charakterizaci cívek a rezonančních obvodů.

\subsection{Ztrátový činitel}

Ztrátový činitel $D$ (dissipation factor) je převrácenou hodnotou činitele kvality:
\begin{equation}
D = \frac{1}{Q} = \frac{R}{|X|} = \frac{\text{Re}(Z)}{|\text{Im}(Z)|}
\end{equation}

Tato veličina se často používá pro charakterizaci kondenzátorů, kde vyjadřuje poměr ztrátové energie k energii uložené v elektrickém poli \cite{KeysightImpedanceHandbook}.

\subsection{Ekvivalentní sériový odpor}

Ekvivalentní sériový odpor \acs{ESR} (Equivalent Series Resistance) představuje rezistivní složku impedance v sériovém náhradním modelu součástky \cite{KeysightImpedanceHandbook}. Pro kondenzátory a cívky je \acs{ESR} důležitým parametrem, který charakterizuje jejich ztráty:
\begin{equation}
\text{ESR} = \text{Re}(Z)
\end{equation}

Nízká hodnota \acs{ESR} je žádoucí zejména u kondenzátorů používaných ve spínaných zdrojích a vysokofrekvenčních aplikacích.


\section{Metody měření impedance}
Existuje mnoho metod pro měření impedance, z nichž každá má své výhody a nevýhody. Neexistuje jediná nejlepší metoda -- volba závisí na požadované přesnosti, měřeném frekvenčním rozsahu, typu měřené impedance a dalších faktorech.

Níže je popsáno několik metod, které se v běžné laboratorní praxi používají. Dále existuje několik specializovaných metod pro měření impedance při rádiových frekvencích, které zde nejsou zmíněny.


\subsection{Můstková metoda}

Měření pomocí můstku je jednoduché a přesné, avšak často nepraktické.
Princip metody spočívá v tom, že pokud jsou impedance obou větví ($Z2$ a $Zx$) stejné,
pak musí být napětí mezi větvemi nulové. Schéma můstku je zobrazeno v obr. \ref{fig:wheatstone}. \cite{KeysightImpedanceHandbook}

\begin{figure}[h!]
    \centering
    \includegraphics[width=0.55\textwidth]{obrazky/wheatstone.pdf}
    \caption{Schéma měřícího můstku}
    \label{fig:wheatstone}
\end{figure}

Když se $Z2$ nastaví tak, aby mezi větvemi bylo naměřeno nulové napětí, pak je měřená
impedance $Zx = \frac{Z1}{Z3}\cdot Z2$

Problém této metody spočívá v tom, že je nutno znát nastavenou hodnotu $Z2$.
Protože nastavitelné kondenzátory a cívky nejsou praktické, používá se tato metoda
zejména pro měření malých změn v odporu, např. tenzometru.

\subsection{Rezonanční metoda}

Rezonanční metoda měření impedance využívá vlastností rezonančního obvodu, který vznikne zapojením měřené impedance do sériového nebo paralelního RLC obvodu. \cite{KeysightImpedanceHandbook}

Princip metody spočívá v tom, že při rezonanční frekvenci $f_0$ dosahuje sériový RLC obvod minimální impedance (pouze rezistivní složka), zatímco paralelní RLC obvod dosahuje maximální impedance. Rezonanční frekvence je dána Thomsonovým vzorcem:
\begin{equation}
f_0 = \frac{1}{2\pi\sqrt{LC}}
\end{equation}

Pro měření neznámé impedance se obvykle používá sériový rezonanční obvod. Měřená impedance se zapojí sériově s opačným typem reaktance známé hodnoty:
\begin{itemize}
    \item Pro měření \textbf{indukčnosti} (cívky) se zapojí sériově s kondenzátorem známé kapacity $C$. Z naměřené rezonanční frekvence $f_0$ se pak vypočítá indukčnost:
    \begin{equation}
    L = \frac{1}{(2\pi f_0)^2 C}
    \end{equation}
    \item Pro měření \textbf{kapacity} (kondenzátoru) se zapojí sériově s cívkou známé indukčnosti $L$. Z naměřené rezonanční frekvence $f_0$ se pak vypočítá kapacita:
    \begin{equation}
    C = \frac{1}{(2\pi f_0)^2 L}
    \end{equation}
\end{itemize}

Výhodou této metody je její jednoduchost a možnost měření při vysokých frekvencích. Nevýhodou je nutnost znalosti přesné hodnoty referenční součástky a omezení na měření jen jedné z reaktivních složek (indukčnosti a kapacity), zatímco rezistivní složka se měří obtížněji.

\subsection{Metoda měření proudu a napětí}
Princip této metody je výpočet hodnoty neznámé impedance z naměřených hodnot napětí a proudu podle Ohmova zákona. \cite{KeysightImpedanceHandbook}
V praxi se proud měří pomocí měření napětí na bočníku.
\begin{equation}
|Z| = \frac{U_{x}}{I_{x}} = \frac{U_{x}}{U_{r}}\cdot R1
\end{equation}
Takovým měřením však získáme pouze modul komplexní impedance, z čehož nelze určit, 
zda je charakter impedance rezistivní nebo reaktivní.
Proto je nutné měřit také fázový posun $U_{x}$ a $U_{r}$, který je roven argumentu $\varphi$.
\begin{equation}
 Z = |Z| \angle \varphi
\end{equation}

\subsection{Samo-vyvažovací můstek}
V praxi se metoda měření proudu a napětí často implementuje pomocí tzv. samo-vyvažovacího můstku,
což je operační zesilovač, který vytváří virtuální zem \vizobr{fig:autobalance_bridge}. \cite{KeysightImpedanceHandbook}

Výhoda takového uspořádání je, že jak měřená impedance, tak bočník jsou připojeny k zemi, 
je tedy možné na nich jednoduše měřit napětí vzhledem k zemi.

\begin{figure}[h!]
    \centering
    \includegraphics[width=0.55\textwidth]{obrazky/i_v_conv.pdf}
    \caption{Schéma samo-vyvažovacího můstku}
    \label{fig:autobalance_bridge}
\end{figure}


\section{Měřiče impedance dostupné na trhu}
Na trhu je mnoho dostupných zařízení pro měření impedance. Pro ukázku byly vybrány tři,
každé reprezentující svou cenovou kategorii.

\subsection{Univerzální tester komponent}

Jako zástupce nejlevnější třídy byl vybrán tester LCR-T4 \vizobr{fig:lcr-t4}. Tento přístroj neměří přímo impedanci,
ale měří hodnoty odporu, kapacity a indukčnosti. Jako doplněk má také funkci rozpoznání tranzistorů.

Pro měření využívá metodu časové konstanty, kdy se měří čas nabití kondenzátoru přes známý odpor.

\begin{figure}[h!]
    \centering
    \includegraphics[width=0.6\textwidth]{obrazky/LCR-T4.jpg}
    \caption{Univerzální tester komponent LCR-T4. Převzato z \cite{hadex_lcrt4}.}
    \label{fig:lcr-t4}
\end{figure}

Přesnost měření není udávána, přístroj totiž slouží pouze pro hrubou identifikaci součástek.
Cena se pohybuje okolo 350~Kč \cite{hadex_lcrt4}.

\subsection{Ruční LCR měřiče}

Pokročilejší přístroje jsou ruční měřiče LCR ve formě podobné multimetrům. Pro příklad
byl vybrán měřič Keysight U1732C \vizobr{fig:keysight_lcr}.

Ruční měřiče jsou vhodné pro rychlou identifikaci cívek a kondenzátorů, nabízejí specifikovanou přesnost
a možnost nastavit parametry měření jako frekvenci měřícího signálu. \cite{Keysight_U1732C_Datasheet}


\begin{figure}[h!]
    \centering
    \includegraphics[width=0.35\textwidth]{obrazky/Keysight_U1732C.jpg}
    \caption{Ruční měřič LCR Keysight U1732C. Převzato z \cite{Keysight_U1732C_Datasheet}.}
    \label{fig:keysight_lcr}
\end{figure}

\subsection{Stolní LCR měřiče}

Pro nejpřesnější měření je nutno použít špičkové stolní LCR měřiče. 
Jako zástupce byl vybrán přístroj LCX200 od firmy Rohde \& Schwarz \vizobr{fig:rhode_lcr}.

Tyto přístroje podporují čtyřbodovou metodu měření, která eliminuje parazitní indukčnost
měřících sond. Na rozdíl od ručních měřičů se zde používají stíněné koaxiální sondy, které jsou méně náchylné k rušení.

Výrobce standardně dodává několik různých fixtur k připojení měřeného vzorku, typicky SMD součástky.

Většina stolních měřičů také podporuje připojení do sítě LAN a automatizaci měření. \cite{RS_LCX_Datasheet}

\begin{figure}[H]
    \centering
    \includegraphics[width=0.85\textwidth]{obrazky/Rhode_schwarz-LCX200.jpg}
    \caption{Stolní měřič Rohde \& Schwarz LCX200. Převzato z \cite{RS_LCX_Datasheet}.}
    \label{fig:rhode_lcr}
\end{figure}

V tabulce \ref{tab:comparison} jsou porovnány přístroje Keysight U1732C a Rohde \& Schwarz LCX200.
Stolní přístroj je výrazně přesnější a nabízí pokročilejší možnosti měření, je ale i násobně dražší.

\begin{table}[h!]
    \centering
    \caption{Porovnání LCR metrů \cite{RS_LCX_Datasheet} \cite{Keysight_U1732C_Datasheet}}
    \label{tab:comparison}
    \vspace{0.2cm}
    
    % Definice sloupců:
    % l = zarovnání doleva (pro první sloupec, aby byl kompaktní)
    % X = sloupec, který vyplní zbývající místo a zalamuje text
    % >{\raggedright\arraybackslash} = zajistí zarovnání vlevo uvnitř X sloupce (lépe vypadá než plné zarovnání do bloku u úzkých sloupců)
    \begin{tabularx}{\textwidth}{@{} l >{\raggedright\arraybackslash}X >{\raggedright\arraybackslash}X @{}}
    \toprule
    \textbf{Parametr}               & \textbf{Keysight U1732C}                  & \textbf{Rohde \& Schwarz LCX200}                  \\ \midrule
    \textbf{Konstrukce}             & Ruční                                     & Stolní                                            \\ \addlinespace
    \textbf{Cena}                   & $\approx 13\,000$ – $15\,000$ Kč              & $\approx 120\,000$ – $300\,000$ Kč                   \\ \addlinespace
    \textbf{Frekvenční rozsah}      & Pevné body: $100$ Hz, $120$Hz, $1$ kHz, $10$ kHz & Spojitý: $4$ Hz až $10$ MHz                       \\ \addlinespace
    \textbf{Základní přesnost}      & $0,2\,\%$                                 & $0,05\,\%$                                        \\ \addlinespace
    \textbf{Amplituda test. signálu}& Pevná (typicky $0,74$ Vrms)               & Nastavitelná: $10$ mV až $10$ V                   \\ \addlinespace
    \textbf{DC Bias (Předpětí)}     & Není k dispozici                          & Interní $0$--$10$ V (Externí až $40$ V)           \\ \addlinespace
    \textbf{Rychlost měření}        & $\sim 1$--$2$ měření/s                    & až $250$ měření/s                   \\ \addlinespace
    \textbf{Konektivita}            & IR-to-USB                                 & USB, LAN, Digital I/O, GPIB                       \\ \bottomrule
    \end{tabularx}
    
    
\end{table}